% ROZDZIAŁ 9 — samodzielny fragment do włączenia w main.tex
\chapter{Rodzaje zbieżności funkcji mierzalnych}
\label{ch:rodzaje-zbieznosci}

\section{Dlaczego potrzebujemy kilku pojęć zbieżności?}

Zbieżność punktowa opisuje zachowanie funkcji w każdym ustalonym
punkcie, lecz nie kontroluje zbiorów, na których funkcje przyjmują
duże wartości. Zbieżność według miary kontroluje rozmiar takich
zbiorów, ale dopuszcza ich przemieszczanie. Zbieżność w całce
kontroluje dodatkowo całkowitą wielkość błędu. To właśnie różnice
między tymi warunkami wyjaśniają, dlaczego nie można dowolnie
zamieniać kolejności granicy i całkowania.

Przez cały rozdział $(X,\mathcal F,\mu)$ jest przestrzenią
z miarą; nie wymagamy zupełności ani $\sigma$-skończoności, chyba
że zaznaczono inaczej. Wszystkie funkcje $f_n,f$ są mierzalne i
rzeczywiste, o ile nie wprowadzono innych założeń. Mówimy, że
własność zachodzi \emph{$\mu$-prawie wszędzie}, gdy zbiór punktów,
w których może zawodzić, zawiera się w pewnym zbiorze należącym
do $\mathcal F$ o mierze zero. Jest to ważne przy miarach
niezupełnych.

\section{Definicje i pierwsze relacje}

Zbieżność według miary zdefiniowaliśmy już w
Definicji~\ref{def:convergence-measure}, a twierdzenie o podciągu
udowodniliśmy w Twierdzeniu~\ref{thm:measure-subsequence}.
Teraz porównujemy ją z pozostałymi pojęciami zbieżności.


\begin{definicja}[Pięć rodzajów zbieżności]
\label{def:five-convergences}
Niech $(X,\mathcal F,\mu)$ będzie przestrzenią z miarą,
a $f_n,f:X\to\mathbb R$ funkcjami mierzalnymi.
\begin{enumerate}[label=\textup{(\roman*)}]
\item $f_n\to f$ \emph{punktowo}, jeśli
$f_n(x)\to f(x)$ dla każdego $x\in X$.
\item $f_n\to f$ \emph{$\mu$-prawie wszędzie}, w skrócie
$f_n\xrightarrow{\mathrm{p.w.}}f$, jeśli zbieżność
punktowa zachodzi poza pewnym mierzalnym zbiorem zerowym.
\item $f_n\to f$ \emph{jednostajnie}, jeśli dla każdego
$\varepsilon>0$ istnieje $N$, takie że
$|f_n(x)-f(x)|<\varepsilon$ dla wszystkich $n\ge N$
i $x\in X$.
\item $f_n\to f$ \emph{prawie jednostajnie}, jeśli dla każdego
$\eta>0$ istnieje $E\in\mathcal F$ z $\mu(E)<\eta$
takie, że $f_n\to f$ jednostajnie na $X\setminus E$.
\item $f_n\to f$ \emph{według miary}, co oznaczamy
$f_n\xrightarrow{\mu}f$, jeśli dla każdego $\varepsilon>0$
zachodzi
\begin{equation}\label{eq:conv-in-measure-9}
\lim_{n\to\infty}\mu\bigl(\{x\in X:
|f_n(x)-f(x)|>\varepsilon\}\bigr)=0.
\end{equation}
\end{enumerate}
Nie wymagamy, aby zbiory poziomicowe w
\eqref{eq:conv-in-measure-9} miały skończoną miarę dla
poszczególnych $n$; wymagana jest ich granica równa zeru.
\end{definicja}

\begin{uwaga}[Dwie kwantyfikacje, których nie wolno mylić]
W zbieżności jednostajnej kolejność kwantyfikatorów jest
następująca: dla każdego $\varepsilon>0$ istnieje jedno $N$
działające \emph{jednocześnie} dla wszystkich $x$. W zbieżności
punktowej indeks $N$ może zależeć od $x$. Zbieżność prawie
jednostajna pozwala usunąć zbiór o dowolnie małej mierze,
ale zbiór ten zależy od wymaganego dopuszczalnego błędu w mierze.
Zbiór wyjątkowy nie może natomiast zależeć od indeksu $n$
w definicji jednostajności na jego dopełnieniu.
\end{uwaga}

\begin{twierdzenie}[Podstawowe implikacje]
\label{thm:conv-implications-9}
Dla dowolnej przestrzeni z miarą $(X,\mathcal F,\mu)$
i mierzalnych funkcji $f_n,f:X\to\mathbb R$ zachodzą:
\begin{enumerate}[label=\textup{(\roman*)}]
\item zbieżność jednostajna pociąga zbieżność punktową
i według miary;
\item zbieżność prawie jednostajna pociąga zbieżność
prawie wszędzie i według miary;
\item zbieżność punktowa pociąga zbieżność prawie wszędzie.
\end{enumerate}
Jeżeli dodatkowo $\mu(X)<\infty$, to zbieżność prawie
wszędzie pociąga zbieżność według miary.
\end{twierdzenie}
\begin{proof}
W (i) dla ustalonego $\varepsilon>0$ i wszystkich dużych
$n$ zbiór $\{|f_n-f|>\varepsilon\}$ jest pusty.
Implikacja punktowa jest oczywista, podobnie (iii).

Dla (ii) ustalmy $\varepsilon,\eta>0$ i wybierzmy
$E\in\mathcal F$ o $\mu(E)<\eta$, poza którym zbieżność
jest jednostajna. Dla dużych $n$ zbiór
$\{|f_n-f|>\varepsilon\}$ zawiera się w $E$,
zatem jego miara jest nie większa od $\eta$.
Ponieważ $\eta$ jest dowolne, otrzymujemy
\eqref{eq:conv-in-measure-9}.

Niech teraz $N$ będzie zbiorem, na którym zbieżność
punktowa nie zachodzi. Ponieważ
\[
N=\bigcup_{m\ge1}\bigcap_{N_0\ge1}
\bigcup_{n\ge N_0}\{|f_n-f|\ge 1/m\},
\]
zbiór $N$ jest mierzalny. Jeśli zbieżność jest prawie
jednostajna, to dla każdego $\eta>0$ istnieje zbiór
$E$ o mierze mniejszej niż $\eta$, poza którym zbieżność
zachodzi; stąd $N\subseteq E$ i $\mu(N)\le\eta$.
Wynika $\mu(N)=0$, co dowodzi drugiej części (ii).

Ostatnia implikacja jest dokładnie twierdzeniem~\ref{thm:ae-to-measure-finite}, już udowodnionym w rozdziale~\ref{ch:funkcje-mierzalne}. To kończy dowód.
\end{proof}

\begin{przyklad}[Dlaczego w ostatniej implikacji potrzebna jest skończoność miary]
\label{ex:ae-not-measure-infinite-9}
Na $\mathbb R$ z miarą Lebesgue'a niech
$f_n=\chi_{[n,\infty)}$. Dla każdego ustalonego
$x\in\mathbb R$ mamy $f_n(x)\to0$, natomiast
dla $0<\varepsilon<1$ zachodzi
$\lambda_1(\{|f_n|>\varepsilon\})=\infty$ dla
każdego $n$. Nie wystarcza nawet zbieżność punktowa
na całej przestrzeni; problemem jest brak kontroli
masy poza zbiorami skończonej miary.
\end{przyklad}

\section{Twierdzenie Jegorowa}

\begin{twierdzenie}[Jegorow]
\label{thm:egorov-9}
Niech $(X,\mathcal F,\mu)$ będzie przestrzenią z miarą
skończoną, tj. $\mu(X)<\infty$. Jeśli
$f_n,f:X\to\mathbb R$ są mierzalne oraz
$f_n\to f$ $\mu$-prawie wszędzie, to
$f_n\to f$ prawie jednostajnie. Dokładniej: dla każdego
$\eta>0$ istnieje $E\in\mathcal F$, $\mu(E)<\eta$,
takie, że
$\sup_{x\in X\setminus E}|f_n(x)-f(x)|\to0$.
Jeśli $X\setminus E=\varnothing$, przyjmujemy supremum
nieujemnej funkcji na zbiorze pustym równe $0$;
zbieżność jednostajna na zbiorze pustym zachodzi.
\end{twierdzenie}
\begin{proof}
Niech $N\in\mathcal F$ będzie zbiorem miary zero,
poza którym $f_n(x)\to f(x)$. Dla $m,N_0\in\mathbb N$
określmy
\[
A_{m,N_0}=\bigcup_{n\ge N_0}
\{x\in X:|f_n(x)-f(x)|\ge1/m\}.
\]
Zbiory $A_{m,N_0}$ są mierzalne i dla ustalonego $m$
maleją wraz z $N_0$. Zbieżność poza $N$ oznacza,
że $\bigcap_{N_0\ge1}A_{m,N_0}\subseteq N$.
Zatem miara tego przecięcia wynosi zero. Skończoność
miary $X$ pozwala zastosować ciągłość z góry:
$\mu(A_{m,N_0})\to0$ dla $N_0\to\infty$.

Ustalmy $\eta>0$. Dla każdego $m$ możemy wybrać
$N_m$ tak duże, by
$\mu(A_{m,N_m})<\eta 2^{-(m+1)}$.
Połóżmy $E=\bigcup_{m\ge1}A_{m,N_m}$.
Przeliczalna podaddytywność daje
$\mu(E)<\eta$.
Dla $x\notin E$ i każdego $m$ oraz $n\ge N_m$
mamy $|f_n(x)-f(x)|<1/m$. Jeśli teraz
$\varepsilon>0$, wybieramy $m$ z $1/m<\varepsilon$;
wówczas nierówność z indeksem $N_m$ zachodzi
jednocześnie dla wszystkich $x\in X\setminus E$.
To jest dokładnie zbieżność jednostajna na
$X\setminus E$.
\end{proof}

\begin{wniosek}[Wersja lokalna dla przestrzeni $\sigma$-skończonej]
\label{cor:local-egorov-9}
Niech $(X,\mathcal F,\mu)$ będzie przestrzenią z miarą,
a $f_n,f:X\to\mathbb R$ funkcjami mierzalnymi.
Załóżmy, że $X=\bigcup_{j\ge1}X_j$, gdzie $X_j\in\mathcal F$,
$X_j\subseteq X_{j+1}$ i $\mu(X_j)<\infty$.
Jeżeli $f_n\to f$ prawie wszędzie, to dla każdego
$\eta>0$ istnieje $E\in\mathcal F$ o $\mu(E)<\eta$
takie, że $f_n\to f$ jednostajnie na
$X_j\setminus E$ dla każdego ustalonego $j$.
Nie twierdzimy, że zbieżność jest jednostajna
na całym $X\setminus E$.
\end{wniosek}
\begin{proof}
Dla każdego $j$ stosujemy twierdzenie
\ref{thm:egorov-9} do ograniczenia miary i funkcji
na $X_j$, z tolerancją $\eta 2^{-(j+1)}$.
Otrzymujemy $E_j\subseteq X_j$ takie, że
$\mu(E_j)<\eta 2^{-(j+1)}$ i zbieżność jest
jednostajna na $X_j\setminus E_j$.
Niech $E=\bigcup_jE_j$. Wtedy $\mu(E)<\eta$, a
$X_j\setminus E\subseteq X_j\setminus E_j$,
więc zbieżność jest jednostajna na pierwszym z tych
zbiorów.
\end{proof}

\begin{przyklad}[Nie istnieje globalna wersja Jegorowa na $\mathbb R$]
Niech $f_n=\chi_{[n,n+1)}$ na $\mathbb R$.
Funkcje zbiegają do zera wszędzie, ale jeśli
$\lambda_1(E)<\infty$, to nie może zachodzić
$f_n\to0$ jednostajnie na $\mathbb R\setminus E$.
W przeciwnym razie istniałoby $N$ takie, że
$[n,n+1)\subseteq E$ dla wszystkich $n\ge N$
(gdyż $f_n$ przyjmuje wyłącznie wartości $0,1$).
Wynikałoby $\lambda_1(E)=\infty$, sprzeczność.
\end{przyklad}

\section{Zbieżność według miary i twierdzenie Riesza o podciągu}

Jednoznaczność granicy według miary, z dokładnością do zbioru zerowego,
została wykazana w Wniosku~\ref{cor:measure-limit-unique}.
Ponadto twierdzenie Riesza o podciągu jest już w pełni
udowodnione w Twierdzeniu~\ref{thm:measure-subsequence};
w dalszych wynikach korzystamy właśnie z tej wcześniejszej wersji.

\begin{wniosek}[Kryterium za pomocą wszystkich podciągów]
\label{cor:every-subseq-9}
Niech $(X,\mathcal F,\mu)$ będzie przestrzenią z miarą,
$\mu(X)<\infty$. Dla mierzalnych
$f_n,f:X\to\mathbb R$ następujące
warunki są równoważne:
\begin{enumerate}[label=\textup{(\roman*)}]
\item $f_n\xrightarrow{\mu}f$;
\item każdy podciąg $(f_{n_k})$ ma dalszy podciąg,
który zbiega do $f$ $\mu$-prawie wszędzie.
\end{enumerate}
\end{wniosek}
\begin{proof}
Jeżeli zachodzi (i), każdy podciąg także zbiega
według miary, więc stosujemy twierdzenie
\ref{thm:measure-subsequence}.
Załóżmy (ii) i przypuśćmy, że (i) nie zachodzi.
Wówczas istnieją $\varepsilon,\delta>0$ oraz
podciąg $(f_{n_k})$ taki, że
$\mu(\{|f_{n_k}-f|>\varepsilon\})\ge\delta$
dla każdego $k$. Na mocy (ii) ma on dalszy
podciąg zbieżny do $f$ prawie wszędzie.
Ponieważ $\mu(X)<\infty$, twierdzenie
\ref{thm:conv-implications-9} daje jego zbieżność
według miary, wbrew ustalonemu oszacowaniu.
\end{proof}

\begin{uwaga}[Nie można skreślić założenia $\mu(X)<\infty$]
Na przestrzeni o mierze nieskończonej warunek (ii)
nie wystarcza: ciąg $\chi_{[n,\infty)}$ na
$\mathbb R$ zbiega wszędzie do zera, lecz nie
według miary. Właśnie dlatego równoważność
\ref{cor:every-subseq-9} została sformułowana
dla miar skończonych.
\end{uwaga}

\section{Zbieżność w $L^p$ i oszacowania Czebyszewa}

\begin{definicja}[Zbieżność w średniej rzędu $p$]
\label{def:lp-convergence-9}
Niech $(X,\mathcal F,\mu)$ będzie przestrzenią z miarą
i niech $1\le p<\infty$. Mierzalną funkcję
$h:X\to\mathbb R$ nazywamy \emph{$p$-całkowalną},
gdy $\int_X|h|^p\,\mathrm d\mu(x)<\infty$.
Dla $p$-całkowalnych funkcji $f_n,f:X\to\mathbb R$
mówimy, że $f_n\to f$ \emph{w $L^p(\mu)$}, jeśli
$\int_X|f_n(x)-f(x)|^p\,\mathrm d\mu(x)\to0$.
Pełną konstrukcję przestrzeni normowanej $L^p$
na klasach funkcji równych prawie wszędzie,
wraz z nierównościami H\"oldera i Minkowskiego,
przedstawimy w rozdziale 14.
\end{definicja}

\begin{twierdzenie}[Nierówność Czebyszewa i zbieżność według miary]
\label{thm:lp-implies-measure-9}
Niech $(X,\mathcal F,\mu)$ będzie przestrzenią z miarą,
$1\le p<\infty$, a $f_n,f:X\to\mathbb R$
funkcjami mierzalnymi. Dla każdego $\varepsilon>0$ zachodzi
\begin{equation}\label{eq:chebyshev-lp-9}
\mu\bigl(\{|f_n-f|>\varepsilon\}\bigr)
\le\varepsilon^{-p}
\int_X|f_n(x)-f(x)|^p\,\mathrm d\mu(x).
\end{equation}
W szczególności zbieżność w $L^p(\mu)$ pociąga
zbieżność według miary.
\end{twierdzenie}
\begin{proof}
Stosujemy nierówność Markowa z rozdziału 7
do nieujemnej mierzalnej funkcji $|f_n-f|^p$
i poziomu $\varepsilon^p$. Jeśli jej całka dąży
do zera, prawa strona \eqref{eq:chebyshev-lp-9}
dąży do zera dla każdego ustalonego $\varepsilon>0$.
\end{proof}

\begin{przyklad}[Zbieżność w $L^p$ nie wymusza zbieżności prawie wszędzie]
\label{ex:typewriter-lp-9}
Dla $k\ge0$ i $0\le j<2^k$ niech
$g_{2^k+j}=\chi_{[j2^{-k},(j+1)2^{-k})}$
na $[0,1)$ z miarą Lebesgue'a. Dla każdego
$1\le p<\infty$ mamy
$\int_{[0,1)}|g_{2^k+j}|^p\,\mathrm d\lambda_1(x)=2^{-k}$,
więc $g_n\to0$ w $L^p$, a także według miary.
Jednak dla każdego $x\in[0,1)$ w każdym bloku o długości
$2^k$ dokładnie jedna funkcja przyjmuje w $x$
wartość $1$, a pozostałe wartość $0$ (dla $k\ge1$).
Stąd $g_n(x)$ nie ma granicy w żadnym punkcie $[0,1)$.
Jest to klasyczna \emph{maszyna do pisania}.
\end{przyklad}

\begin{przyklad}[Zbieżność prawie wszędzie i według miary bez zbieżności w $L^1$]
Niech $f_n=n\chi_{(0,1/n)}$ na $(0,1)$.
Wtedy $f_n\to0$ wszędzie i według miary, lecz
$\int_{(0,1)}|f_n|\,\mathrm d\lambda_1(x)=1$
dla każdego $n$. Brakuje jednostajnej całkowalności;
por. twierdzenie Vitaliego z rozdziału 8.
\end{przyklad}

\section{Ciągi Cauchy'ego według miary i zupełność}

\begin{definicja}[Ciąg Cauchy'ego według miary]
\label{def:cauchy-measure-9}
Niech $(X,\mathcal F,\mu)$ będzie przestrzenią z miarą.
Ciąg mierzalnych funkcji rzeczywistych
$f_n:X\to\mathbb R$ nazywamy \emph{ciągiem Cauchy'ego według miary}, jeśli
dla wszystkich $\varepsilon,\delta>0$ istnieje
$N\in\mathbb N$ takie, że dla $m,n\ge N$
zachodzi $\mu(\{|f_n-f_m|>\varepsilon\})<\delta$.
\end{definicja}

\begin{twierdzenie}[Zupełność zbieżności według miary]
\label{thm:cauchy-measure-complete-9}
Niech $(X,\mathcal F,\mu)$ będzie dowolną przestrzenią
z miarą. Jeżeli $(f_n)$ jest ciągiem Cauchy'ego według
miary złożonym z mierzalnych funkcji rzeczywistych,
to istnieje mierzalna funkcja $f:X\to\mathbb R$ taka,
że $f_n\xrightarrow{\mu}f$. Funkcja $f$ jest wyznaczona
jednoznacznie z dokładnością do równości prawie wszędzie.
\end{twierdzenie}
\begin{proof}
Wybierzmy ściśle rosnący ciąg indeksów $n_k$
rekurencyjnie tak, aby dla wszystkich $m,n\ge n_k$
zachodziło
$\mu(\{|f_m-f_n|>2^{-k}\})<2^{-k}$.
Możemy jednocześnie wymagać $n_{k+1}>n_k$.
Niech
$A_k=\{|f_{n_{k+1}}-f_{n_k}|>2^{-k}\}$.
Wówczas $\sum_k\mu(A_k)<\infty$ i lemat
Borela--Cantellego zapewnia, że poza mierzalnym
zbiorem zerowym $N$ każdy punkt należy tylko do
skończenie wielu $A_k$.

Dla $x\notin N$ i dużych $k$ mamy
$|f_{n_{k+1}}(x)-f_{n_k}(x)|\le2^{-k}$.
Dla $j>k$ nierówność trójkąta daje
\[
|f_{n_j}(x)-f_{n_k}(x)|
\le\sum_{\ell=k}^{j-1}2^{-\ell}\le2^{1-k}.
\]
Zatem $(f_{n_k}(x))$ jest ciągiem Cauchy'ego w
$\mathbb R$ i zbiega do pewnej skończonej liczby
$f(x)$. Na $N$ definiujemy $f(x)=0$.
Mierzalność $f$ wynika z twierdzenia o mierzalności
granic punktowych po przedłużeniu funkcji granicznej
zerem na $N$. Zbieżność $f_{n_k}\to f$ prawie
wszędzie nie gwarantuje sama przez się zbieżności
według miary na przestrzeni nieskończonej.

Tutaj jednak mamy silniejsze oszacowanie: dla
ustalonego $k$ i dowolnego $j>k$, poza zbiorem
$\bigcup_{\ell\ge k}A_\ell$, zachodzi
$|f_{n_j}-f_{n_k}|\le2^{1-k}$.
Przechodząc do granicy $j\to\infty$ poza $N$,
dostajemy
\[
\mu\bigl(\{|f-f_{n_k}|>2^{1-k}\}\bigr)
\le\sum_{\ell\ge k}\mu(A_\ell)
\le\sum_{\ell\ge k}2^{-\ell}\longrightarrow0.
\]
Zatem $f_{n_k}\to f$ według miary.
Dla ustalonych $\varepsilon,\delta>0$ warunek
Cauchy'ego daje $N_C$ takie, że dla $m,n\ge N_C$
zachodzi
$\mu(\{|f_n-f_m|>\varepsilon/2\})<\delta/2$.
Wybieramy teraz $k$ \emph{jednocześnie} tak duże,
aby $n_k\ge N_C$ i
$\mu(\{|f_{n_k}-f|>\varepsilon/2\})<\delta/2$.
Wówczas dla każdego $n\ge N_C$ mamy
$\mu(\{|f_n-f_{n_k}|>\varepsilon/2\})<\delta/2$.
Wówczas $\mu(\{|f_n-f|>\varepsilon\})<\delta$.
Jednoznaczność wynika z twierdzenia
\ref{cor:measure-limit-unique}.
\end{proof}

\section{Metryzacja zbieżności według miary na przestrzeni skończonej miary}

\begin{twierdzenie}[Metryka zbieżności według miary]
\label{thm:measure-metric-9}
Niech $(X,\mathcal F,\mu)$ będzie przestrzenią z miarą,
$\mu(X)<\infty$. Oznaczmy przez
$\mathcal M(X,\mathcal F)$ zbiór mierzalnych funkcji
rzeczywistych na $X$, a przez $\mathcal M/\!\sim$
jego iloraz przez relację
$f\sim g\Longleftrightarrow f=g$ $\mu$-prawie wszędzie.
Dla klas $[f],[g]\in\mathcal M/\!\sim$ połóżmy
\begin{equation}\label{eq:measure-metric-9}
\rho([f],[g])=
\int_X\min\{1,|f(x)-g(x)|\}\,\mathrm d\mu(x).
\end{equation}
Wówczas $\rho$ jest metryką na $\mathcal M/\!\sim$.
Ponadto $f_n\xrightarrow{\mu}f$ wtedy i tylko wtedy,
gdy $\rho([f_n],[f])\to0$.
\end{twierdzenie}
\begin{proof}
Całka w \eqref{eq:measure-metric-9} jest skończona,
bo całkujemy funkcję ograniczoną przez $1$,
a $\mu(X)<\infty$. Równość funkcji prawie
wszędzie nie zmienia wartości całki, więc wzór
zależy tylko od klas. Symetria i nieujemność
są oczywiste. Jeśli $\rho([f],[g])=0$, to
nieujemna funkcja $\min\{1,|f-g|\}$ ma całkę
zero, a więc znika prawie wszędzie; wynika
$[f]=[g]$. Nierówność trójkąta jest konsekwencją
punktowej nierówności
\[
\min\{1,|f-h|\}\le
\min\{1,|f-g|\}+\min\{1,|g-h|\},
\]
którą otrzymujemy z nierówności trójkąta dla
wartości bezwzględnej i oszacowania
$\min\{1,a+b\}\le\min\{1,a\}+\min\{1,b\}$
dla $a,b\ge0$. Całkujemy obie strony.

Załóżmy $\rho([f_n],[f])\to0$. Dla $\varepsilon>0$
zachodzi punktowo
$\min\{1,|f_n-f|\}\ge
\min\{1,\varepsilon\}\chi_{\{|f_n-f|>\varepsilon\}}$.
Zatem
\[
\mu(\{|f_n-f|>\varepsilon\})\le
\frac{\rho([f_n],[f])}{\min\{1,\varepsilon\}}
\longrightarrow0.
\]

Odwrotnie, jeśli $f_n\to f$ według miary,
to dla każdego $0<\varepsilon<1$ mamy
\[
\rho([f_n],[f])\le\varepsilon\mu(X)+
\mu(\{|f_n-f|>\varepsilon\}).
\]
Dla ustalonego $\varepsilon$ drugi składnik
dąży do zera, zatem
$\limsup_n\rho([f_n],[f])\le\varepsilon\mu(X)$.
Pozwalając $\varepsilon\downarrow0$, otrzymujemy
$\rho([f_n],[f])\to0$.
\end{proof}

\begin{wniosek}[Zupełność metryki zbieżności według miary]
\label{cor:measure-metric-complete-9}
Niech $(X,\mathcal F,\mu)$ będzie przestrzenią z miarą,
$\mu(X)<\infty$, a $(\mathcal M/\!\sim,\rho)$ przestrzenią
metryczną klas mierzalnych funkcji rzeczywistych z metryką
zdefiniowaną w twierdzeniu \ref{thm:measure-metric-9}.
Wówczas $(\mathcal M/\!\sim,\rho)$ jest zupełna.
\end{wniosek}
\begin{proof}
Niech $([f_n])$ będzie ciągiem Cauchy'ego w metryce
$\rho$. Dla $\varepsilon>0$ mamy
$\mu(\{|f_n-f_m|>\varepsilon\})\le
\rho([f_n],[f_m])/\min\{1,\varepsilon\}$.
Zatem $(f_n)$ jest ciągiem Cauchy'ego według miary.
Z twierdzenia~\ref{thm:cauchy-measure-complete-9}
istnieje mierzalna $f$ taka, że $f_n\to f$
według miary. Twierdzenie~\ref{thm:measure-metric-9}
daje $\rho([f_n],[f])\to0$.
\end{proof}

\begin{uwaga}[Skończoność miary jest istotna dla tego wzoru]
Jeśli $\mu(X)=\infty$, wartość
$\int_X\min\{1,|f-g|\}\,\mathrm d\mu(x)$
może być nieskończona, więc nie musi określać
metryki o wartościach skończonych.
Na przestrzeniach $\sigma$-skończonych można budować
metryki z wagami po wyczerpaniu przestrzeni zbiorami
o skończonej mierze, ale taka metryka na ogół
opisuje zbieżność \emph{lokalną} według miary,
nie globalną zbieżność w znaczeniu
\eqref{eq:conv-in-measure-9}.
\end{uwaga}


\begin{przyklad}[Zbieżność lokalna według miary nie jest globalną]
\label{ex:local-vs-global-9}
Na $\mathbb R$ z miarą Lebesgue'a rozważmy
$f_n=\chi_{[n,n+1)}$ i $f=0$. Dla każdego
$R>0$ mamy $f_n|_{[-R,R]}=0$ dla dostatecznie
dużych $n$; funkcje zbiegają zatem według miary na
każdym zbiorze $[-R,R]$. Jednak dla $0<\varepsilon<1$
zachodzi $\lambda_1(\{|f_n-f|>\varepsilon\})=1$
dla wszystkich $n$. Nie ma globalnej zbieżności
według miary. Ten przykład objaśnia różnicę między
metryzacją zbieżności lokalnej na przestrzeni
$\sigma$-skończonej i metryką \eqref{eq:measure-metric-9}
na przestrzeni o mierze skończonej.
\end{przyklad}

\begin{twierdzenie}[Kontrola ucieczki masy przy zbieżności lokalnej]
\label{thm:local-to-global-measure-9}
Niech $(X,\mathcal F,\mu)$ będzie przestrzenią z miarą,
a $(X_j)_{j\ge1}$ ciągiem zbiorów mierzalnych.
Załóżmy, że $X_j\in\mathcal F$, $X_j\subseteq X_{j+1}$,
$\bigcup_jX_j=X$ i $\mu(X_j)<\infty$.
Niech $f_n,f:X\to\mathbb R$ będą mierzalne.
Wówczas $f_n\to f$ według miary na całym $X$
wtedy i tylko wtedy, gdy dla każdego $\varepsilon>0$:
\begin{enumerate}[label=\textup{(\roman*)}]
\item $\mu(X_j\cap\{|f_n-f|>\varepsilon\})\to0$
dla każdego ustalonego $j$;
\item $\displaystyle \lim_{j\to\infty}\limsup_{n\to\infty}
\mu((X\setminus X_j)\cap\{|f_n-f|>\varepsilon\})=0$.
\end{enumerate}
\end{twierdzenie}
\begin{proof}
Oznaczmy $B_n=\{|f_n-f|>\varepsilon\}$.
Jeśli $\mu(B_n)\to0$, oba warunki są natychmiastowe.
Odwrotnie, dla każdego $j,n$ z podaddytywności
\[
\mu(B_n)\le\mu(X_j\cap B_n)+
\mu((X\setminus X_j)\cap B_n).
\]
Przechodząc do $\limsup$ po $n$ i korzystając
z (i), dostajemy
$\limsup_n\mu(B_n)\le
\limsup_n\mu((X\setminus X_j)\cap B_n)$.
Na mocy (ii) prawa strona dąży do zera wraz
z $j\to\infty$, co dowodzi tezy.
\end{proof}

\section{Użyteczne kryteria graniczne}

\begin{twierdzenie}[Lemat Scheffégo]
\label{thm:scheffe-9}
Niech $(X,\mathcal F,\mu)$ będzie przestrzenią z miarą,
a $f_n,f:X\to[0,\infty)$ funkcjami mierzalnymi,
całkowalne oraz $f_n\to f$ prawie wszędzie.
Jeżeli
$\int_Xf_n(x)\,\mathrm d\mu(x)\to
\int_Xf(x)\,\mathrm d\mu(x)<\infty$,
to
$\int_X|f_n(x)-f(x)|\,\mathrm d\mu(x)\to0$.
\end{twierdzenie}
\begin{proof}
Mamy punktową tożsamość
$|f_n-f|=f_n+f-2\min(f_n,f)$.
Ponieważ $0\le\min(f_n,f)\le f$ oraz
$\min(f_n,f)\to f$ prawie wszędzie,
twierdzenie o zbieżności zmajoryzowanej daje
$\int_X\min(f_n,f)\,\mathrm d\mu(x)\to\int_Xf\,\mathrm d\mu(x)$.
Całkując tożsamość i korzystając z założenia
o granicach całek, otrzymujemy tezę.
\end{proof}

\begin{wniosek}[Wersja dla gęstości prawdopodobieństwa]\label{idx:09:wniosek:14}
Niech $(X,\mathcal F,\mu)$ będzie przestrzenią z miarą,
a $f_n,f:X\to[0,\infty)$ funkcjami mierzalnymi.
Jeżeli
$\int_X f_n\,\mathrm d\mu(x)=
\int_X f\,\mathrm d\mu(x)=1$ oraz $f_n\to f$ prawie
wszędzie, to $\int_X|f_n-f|\,\mathrm d\mu(x)\to0$.
W zastosowaniach do prawdopodobieństwa jest to
podstawowy sposób przejścia od zbieżności gęstości
prawie wszędzie do zbieżności w wariacji całkowitej.
\end{wniosek}

\begin{uwaga}[Co naprawdę wynika z poszczególnych pojęć?]
\label{rem:conv-summary-9}
Na przestrzeni skończonej miary zbieżność prawie
wszędzie pociąga zbieżność według miary; na dowolnej
przestrzeni zbieżność według miary zapewnia podciąg
zbieżny prawie wszędzie. Ani zbieżność według miary,
ani w $L^p$ nie muszą zapewniać zbieżności prawie
wszędzie całego ciągu. Sama zbieżność prawie wszędzie
nie zapewnia zbieżności całek. Dodatkową kontrolę
zapewniają dominacja, monotoniczność lub jednostajna
całkowalność.
\end{uwaga}

\begin{center}
\begin{tikzpicture}[>=stealth, every node/.style={font=\small},
box/.style={draw,rounded corners,align=center,inner sep=5pt}]
\node[box] (u) at (0,2.5) {jednostajnie};
\node[box] (au) at (4.4,2.5) {prawie jednostajnie};
\node[box] (ae) at (4.4,0.7) {prawie wszędzie};
\node[box] (lp) at (-0.1,-1) {w $L^p$, $1\le p<\infty$};
\node[box] (mu) at (4.4,-1) {według miary};
\draw[->] (u) -- (au);
\draw[->] (au) -- (ae);
\draw[->] (au.east) to[out=0,in=0] (mu.east);
\draw[->] (lp) -- (mu);
\draw[->,dashed] (ae) -- node[right,font=\scriptsize]{gdy $\mu(X)<\infty$} (mu);
\draw[->,dashed,bend left=24] (ae.west) to node[above,font=\scriptsize]{Jegorow: $\mu(X)<\infty$} (au.south west);
\end{tikzpicture}
\par\smallskip\textit{Schemat implikacji. Strzałki przerywane wymagają skończoności miary.}
\end{center}
